%% kapa: exam=1 round=32 year=2021 period=1 form=A subject=real_estate subject_ko=부동산학원론 seq=1 number=81
%% source: Q-Net 감정평가사 기출문제 (cst00309) artlSeq=5208397 file=extracted/1008/5208397/A형/2021년도 제32회 감정평가사 1차 1교시 A형.pdf
%% conversion: kordoc 4.24.51 -> tools/gen1.py (auto); review=proofread-sonnet; flags=none
\begin{question}{81}
부동산의 개념에 관한 설명으로 옳지 않은 것은?
\begin{choices}
  \item 자연·공간·위치·환경 속성은 물리적 개념에 해당한다.
  \item 부동산의 절대적 위치는 토지의 부동성에서 비롯된다.
  \item 토지는 생산의 기본요소이면서 소비재가 된다.
  \item 협의의 부동산과 준부동산을 합쳐 광의의 부동산이라고 한다.
  \item 부동산의 법률적·경제적·물리적 측면을 결합한 개념을 복합부동산이라고 한다.
\end{choices}
\end{question}
\answer{5}
